\section{Marked disk laws and uniform boundary control}\label{sec:proxy}

We need boundary control that survives the limit $\gamma\uparrow2$ and
applies to the actual maps of cut disks into their parent.  We establish
the marked disk laws and their boundary measures first.  A fixed strictly
subcritical metric will then convert small quantum boundary arcs and
small offspring perimeters into small images in the parent chart.

\subsection{The disk normalization and strip representation}\label{proxy:disk-normalization}

We use \(\kappa=\gamma^2\in(8/3,4)\), increasing to 4 (thus using \emph{simple}, not FK-peanosphere, subcritical explorations in this argument), \(r=4/\kappa\) and \(a=Q_\gamma-\gamma\) as in Section~\ref{sec:continuum}.
All disk areas in the fields considered here (including the sphere fields in Equation~\eqref{eq:original-1}) are in units \(\mu^\gamma/(2(2-\gamma))\).
For boundary length start with the usual flat radius-normalized free-boundary chaos with parameter \(\gamma/2\), divided by \(2-\gamma\); we will multiply this by a deterministic scalar so that a DMS quantum disk at perimeter \(s\) has area density
\[
 d_s^r(v)=\frac{s^{2r}v^{-1-r}}{\Gamma(r)}e^{-s^2/v},\qquad v>0 .
\]
Here we use the Ang--Gwynne quantum disk area law \citep[Theorem~1.2]{AngGwynneDisk} (ordinary \(\gamma\)-LQG disk of given boundary length, \(\gamma\in(0,2)\)), and homogeneity.
The scalar converges to a strictly positive finite limit as shown below.
At criticality we take these limiting units.
They agree for area with the derivative convention used above.
As throughout, assertions for \(\gamma<2\) in limiting estimates need hold only for \(\gamma\) sufficiently close to 2.

On the strip \({\cal S}=\mathbb R\times(0,\pi)\) the doubly marked disk has the analogue of \eqref{eq:original-1}.
Replace \(G^\perp\) by the lateral Neumann GFF \(P^\perp\) (the nonconstant cosine modes) and \(R_a\) by the following process, where $W^+$ and $W^-$ are independent standard three-dimensional Brownian motions:
\[
 U_a(t)=|W^{\operatorname{sgn}(t)}_{2|t|}+a |t|e_1|.
\]
Write \(V_\gamma\) for the total initial-units normalized boundary length (division by \(2-\gamma\) but not yet the extra scalar) of \(P^\perp-U_a\).
Given initial-units perimeter 1 the field is \(P^\perp-U_a-(2/\gamma)\log V_\gamma\), with tilt proportional to \(V_\gamma^{2a/\gamma}\).
Indeed DMS here use a Bessel excursion of dimension \(3-4/\gamma^2<2\) with the \(2/\gamma\)-log run at quadratic variation 2.
The maximum intensity after taking this log is proportional to \(e^{-a m}dm\), and the same maximum and reversal argument as for \eqref{eq:original-1} gives the shape just stated. DMS boundary marking gives conditionally independent length samples at the two ends given the unmarked length-conditioned quantum disk.
In particular one can also use the disk with just one such typical end, add the other independently by length to give strip coordinates, and locate any further \emph{specified relative length labels} instead of independent marks.
Given the two end marks split at prescribed lengths one may similarly choose strip coordinates using the first mark and an added independent length sample, ignoring the second mark for shape estimates.
Here prescribing the split from a uniform second mark only assigns it a deterministic position in length coordinates, without biasing the once marked disk \citep{DMS,AngGwynneDisk}.

\subsection{Boundary chaos and mass at the strip edges}\label{proxy:chaos}

We now establish the convergence and uniform bounds required below.
For chaos on the free boundary we use the general log-correlated subcritical-to-critical theorem of Powell \citep[Theorem~3.1]{PowellReview}, together with critical-chaos existence, uniqueness and positivity \citep[Theorems~5.3--5.4]{JSWchaos}.  These results apply locally to the boundary restriction of a Neumann GFF.
Indeed, one can work first in straight coordinates with a mixed GFF (Neumann on the diameter of a disk, Dirichlet on the remaining semicircle), coupled by even reflection to a Dirichlet GFF on the disk.
On compact subintervals of the diameter its covariance divided by $2$ is logarithmic plus a smooth term.  Equivalently, divide the boundary field by $\sqrt2$ and apply the one-dimensional theorem with chaos parameter $\gamma/\sqrt2\uparrow\sqrt2$.
Semicircle/radius normalization in the subcritical chaos agrees by subcritical uniqueness with smooth mollification and Wick normalization up to the usual smooth regular-part factor, convergent to a positive smooth factor; the semicircle centered on the line for the reflected construction uses exactly the full circle-average regular part.
Thus in our initial units we get local convergence in probability to positive diffuse measures of full support (a fixed positive multiple of radius-normalized critical boundary chaos).
Harmonic or continuous random field additions multiply the positive measures as in the interior.
For \(P^\perp\) on a unit strip or slab we may add an independent average Brownian motion to make the free GFF, decompose by the Markov property in a larger patch up to smooth harmonic parts on the observed closure, then undo the added average.
The resulting continuous corrections on a fixed size patch have Gaussian supremum tails.
All estimates obtained patchwise in this way are stationary for \(P^\perp\).
The local area statements away from the boundary use just the interior chaos theorems already applied above \citep{PowellReview,JSWchaos}.

\begin{lemma}[Boundary mass moments]\label{proxy:boundary-moments}
The normalized boundary masses on fixed intervals have uniformly bounded small positive moments and their inverses on nonempty open boundary intervals have all positive moments uniformly. \end{lemma}\begin{proof}
It is enough to prove this for mixed half-disks on a compact subinterval of the flat edge with nonempty interior; continuous Gaussian-tail corrections and Hölder handle the local applications.
For positive moments pack \(c k\) disjoint translated scaled copies of the domain along the subinterval with scale \(c'/k\) using the bulk packing argument of Section~\ref{sec:continuum}.
The Markov complementary harmonic field is independent of the small fields; its restrictions to the copies, estimated on the observed subintervals, have variances \(O(\log k)\) there and on fixed relative enlargements.
The probability of a minimum below \(-B\log k\) is \(O(k^{1-cB^2})\) by reflected harmonic interior estimates. (We only claim independence of the whole complementary field from the small GFFs, not independence between its restrictions.) Tightness of the mass on the original interval follows by the local limit.
Thus exactly as before the maximum of \(c k\) iid masses gives a polynomial tail estimate for the single mass.
Conversely absence of one of these masses \(\ge\delta>0\) has exponentially small probability by uniform positivity.
On the good harmonic event this bounds the original mass below by \(c'' k^{-C(1+B)}\); taking \(B\) large proves every desired inverse moment.
Constants and laws on all rescaled copies use the original standard flat-coordinate normalization.
We will use such inverse estimates also independently of a semicircle coarse tilt: at scale \(e^{-t}\) about a boundary point \(x\) of the Neumann half-plane, write \(H_t\) for the semicircle average at that radius (fix the additive constant by \(H_0=0\) at some given center, or subtract the local starting average).
On tilting by \(\exp(H_t-\operatorname{Var}(H_t)/2)\), the normalized length of an interval at distance and with size comparable to \(e^{-t}\) from \(x\), divided by \(\exp((\gamma/2)(H_t-Q_\gamma t))\), still has inverse moments of all orders uniformly.
This follows by rescaling: the tilt adds a covariance which, after subtracting its semicircle average shift, is bounded on the fixed rescaled interval, by the averaged logarithmic kernel; the residual random field there is a standard rescaled free field with tight Gaussian choice of its residual constant.
The statement is unchanged with tilt by a bounded multiple of \(H_t\) with its variance compensation.
\end{proof}

\begin{lemma}[An interior moment above one half]\label{proxy:bulk-moment}
Normalized near-critical bulk masses of a Dirichlet GFF on a fixed compact disk in the interior have a uniformly bounded moment of some order \(p>1/2\). \end{lemma}\begin{proof}
 In fixed interior nested disks pack \(\asymp k^2\) disks on a mesh of order \(1/k\), of disjoint larger closure, and test copies of a fixed scaled compact inner-disk mass.
The harmonic outside fields \(Y_i\) at centers have variances \(v_i=\log k+O(1)\), with pairwise covariances bounded above by \(\log |x_i-x_j|^{-1}+O(1)\).
Oscillations on the inner disks about the centers have uniform Gaussian tails, also when individually tilted by \(\exp(\alpha Y_i-\alpha^2 v_i/2)\) for fixed \(\alpha\); this uses the Green kernels of the disk and subdisks and harmonic interior bounds.
Choose small \(\alpha>0,\ \alpha^2<2\).
The sum of these weights divided by \(k^2\) has bounded second moment (mesh integral of the covariance exponential).
Keeping only terms with \(|Y_i-\alpha\log k|\le\epsilon\log k\) and oscillation at most \(\epsilon\log k\) its first moment stays bounded below.
Thus with positive probability independently of \(k\) there are at least \(c k^{2-\alpha^2/2-\alpha\epsilon}\) good indices, determined independently of the small GFFs.
On them the scale factor for the small normalized mass to the original mass is at least \(c' k^{-4+2\alpha-C\epsilon}\) uniformly for \(\gamma\) sufficiently close to 2.
Tightness of the original mass and the iid copies now bound the single tail at \(C' k^{4-2\alpha+C\epsilon}\) by \(O(k^{-2+\alpha^2/2+\alpha\epsilon})\), as the above good-count probability stays bounded below.
For small \(\epsilon\) this proves the moment assertion.
\end{proof}

\begin{lemma}[Uniform Gaussian barrier at the boundary]\label{proxy:barrier}
Let $H_n(x)$ be the semicircle average at radius $e^{-n}$ of the locally normalized Neumann field on a fixed compact boundary interval.  Put $X_n(x)=2n-H_n(x)$.  There are uniform constants $C$ such that, for $y\geq0$,
\[
 \Prob\{X_n(x)<-y\text{ for some }x\text{ in the interval and }n\geq1\}
 \leq C(1+y)^C e^{-y}.
\]
The endpoint-tilted walk and the spatial oscillations satisfy the ballot and Gaussian bounds established in the proof.
\end{lemma}\begin{proof}

Consider one unit length of straight boundary in the half-plane, extending tangentially as necessary, with Neumann GFF modulo constants.
Use the free-field representative with reference semicircle average zero of a large radius \(R\) at the midpoint, strictly enclosing the test neighborhoods.
Then \(H_n(x)\), the semicircle averages at radii \(e^{-n}\), are Gaussian with independent Brownian increments of variance rate 2 at each fixed center and uniformly bounded starting variances.
This follows from even reflection of the log kernel; in the test neighborhoods the reference-circle logarithmic potential is constant.
We will use two elementary uniform barrier estimates.
Tilt at \(n\) by \(\exp(H_n(x)-\operatorname{Var}(H_n(x))/2)\); then \(X_i=2i-H_i(x)\) for \(1\le i\le n\) are a centered-diffusion variance-2 Gaussian walk with bounded initial shift law.
Uniformly,
\[
 {\mathbb P}_*(X_n\in[-y+j,-y+j+1],\ X_i\ge-y,\ 1\le i<n)
 \le C(1+y+|j|+\log(2+n))^C n^{-3/2},\quad y\ge0.
\]
This is just the killed Gaussian bridge estimate: between integer times, a Brownian bridge with these endpoints above \(-y\) (except possibly at \(n\)) stays above \(-y-|j|-2-\sqrt{C'\log(2+n)}\) with probability bounded below; then reflection on the whole interval and integration of the initial Gaussian give the bound.
The oscillation \(O_n\) of \(H_n\) over center displacements \(\le C'e^{-n}\) from \(x\) has uniform Gaussian tails, including under the tilt.
Indeed rescaled log-circle-average fields have increment variances bounded by \(C|\Delta x|^\eta\) for some \(\eta>0\) (average the log potential \(-\log(1\vee|z-u|)\)); the tilt shift in increments is bounded.

Except for probability \(C(1+y)^C e^{-y}\), \(X_n(x)\ge-y\) simultaneously on our boundary interval for every \(n\ge1\).
At a global first failing level take a nearest grid center of spacing \(\le e^{-n}\) on the interval; it has survived on all earlier levels and in the endpoint band \(j\) failure needs \(O_n\ge j\).
Undo the tilt and sum over \(O(e^n)\) centers at cost \(C e^{-y+j}\) times the tilted probabilities.
For \(j>0\) or when needed below, Hölder keeps a power \(b>2/3\) of the ballot bound, while retaining a Gaussian oscillation tail; now sum over \(j\in\mathbb Z,n\ge1\).
Compact horizontal extensions introduce no change.

\end{proof}

\begin{lemma}[Full disk convergence]\label{proxy:disk-convergence}
As $\gamma\uparrow2$, the perimeter-conditioned strip fields converge locally as distributions, jointly with their finite area and boundary measures and their prescribed relative boundary-length labels.  The measure and label convergence includes the strip ends.  The deterministic conversion from the initial boundary units to the units in $d_s^r$ converges to a finite positive constant.  In the limiting units the ordinary disk of perimeter $s$ has area density $d_s$ from Proposition~\ref{disk:palm}.  The limiting measures have no boundary or end atoms.
\end{lemma}\begin{proof}

Test the area in interior Whitney cells near this line at layer \(n\) (distance \(\asymp e^{-n}\)) with \(O(e^n)\) cells, of sufficiently small fixed relative diameter.
Choose associated semicircle center on the line and average at a comparable outer scale (e.g. \(n\) replaced by \(n-C_0\) for all cells, an integer).
The cell is in a larger interior disk disjoint from the semicircles at these and all earlier levels.
Up to bounded factors its normalized area is bounded by
\[
 \exp(\gamma(H_{n-C_0}-Q_\gamma n))\,Y\,Z ,
\]
where \(Z\) has the independent rescaled compact Dirichlet mass law and \(Y\) has all moments uniformly under the tilt at this last level (supremum of the complementary field minus \(H_{n-C_0}\), by rescaled Green estimates); all barrier and endpoint tests at this cell center are independent of the Dirichlet field.
Constants can depend on \(C_0\).
Choose \(p\in(1/2,1)\) from the moment just proved, and keep \(\gamma p>1\) uniformly.
On the global good-barrier event, bound the tilted expectation of the \(p\)-power of the cell mass bandwise using independence for \(Z\), Hölder with all moments for \(Y\), and power \(b\in(2/3,1)\) of the ballot bound at \(j\ge0\).
After summing centers this gives at each large \(n\) an upper bound
\[
 C e^{(\gamma p-1)y}(1+y+\log(2+n))^C n^{-3b/2}.
\]
Here \(Q_\gamma\ge2\); no strictly positive charge gap is used.
Thus at fixed \(y\), tails over Whitney layers vanish in probability on the good event.
Taking thresholds of order \(e^{\gamma y}\) also gives uniform small positive moments of total mass by subadditivity, including the finitely many deeper-interior coverings.
This proves the desired estimates for the lateral strip field as well by local Neumann decomposition/comparison (adding a continuous Gaussian-tail term, Hölder); it proves convergence of masses to the interior limit without boundary atoms.

Over the infinite strip, independence and \(\mathbb E\exp(-b\inf_{[k,k+1]}U_a)=O_b(k^{-3/2})\), as in Section~\ref{sec:continuum}, now finish the joint convergence of finite measures, fields locally, and boundary labels up to the ends.
The tilt \(V_\gamma^{2a/\gamma}\) converges to 1 in \(L^1\).
Disk area at fixed initial-units perimeter therefore has a strictly positive finite limiting distribution.
For clarity, the convergence of the boundary conversion scalar follows
already from the area distributions.  If $b_\gamma$ is that scalar, a disk
of perimeter one in the initial units has area law
$b_\gamma^2/E_r$, where $E_r$ has the Gamma$(r,1)$ law.  The area variables
just proved to converge have a finite strictly positive limiting law, while
$E_r$ converges to an exponential variable as $r\downarrow1$.  These facts
first bound $b_\gamma$ away from zero and infinity.  Every subsequential
limit $b$ then gives the law $b^2/E$; its strictly monotone distribution
function, or any one fixed quantile, determines $b$ uniquely.  Thus
$b_\gamma$ converges to a finite positive constant.
In the limiting normalization, given perimeter \(s\), unbiased disk area has density \(d_s=d_s^1\) as in \eqref{eq:original-3}.
The end-mark symmetry passes to criticality; alternatively re-embed using three sampled boundary points by continuity and apply the resampling test as for the spheres in Section~\ref{sec:continuum}.
These disk laws and estimates make no assertion of subcritical LQG \emph{metric} convergence.

\end{proof}

\subsection{Arc measures in random conformal coordinates}\label{proxy:arc-measures}

\begin{lemma}[Agreement on common collars]\label{proxy:arc-covariance}
Boundary chaos can be realized as an intrinsic measurable function of the abstract disk.  On a common open boundary collar, two qualifying disk or wedge charts related by the charge-$Q_\gamma$ coordinate rule give the same boundary measure on that side.  This holds for random analytic coordinate changes and for the critical limiting disk law.  Bulk measures agree on common conformal neighborhoods as well.
\end{lemma}\begin{proof}

\proofstage{Approximation in standard charts}

In disk or strip charts as above the boundary chaos is intrinsic and is determined by the field as an abstract disk (with or without the marks), and on genuinely shared open boundary segments with the same side and conformal collar it agrees between charts.
We need this also for the disk laws occurring as \emph{marginals} in cuttings, with coordinate identifications not necessarily independent of the fields.
Work first in straight local half-plane charts off the distinguished ends, and at a strictly subcritical test exponent \(\eta/2\).
Positive length estimates/measures here refer to \(\epsilon^{\eta^2/4}\exp(\eta h_\epsilon/2)dx\) (deterministic normalizing factors suppressed), and may equivalently be obtained (up to a fixed kernel constant) by smooth approximate averages centered a distance \(\asymp\epsilon\) \emph{inside} the chart, with support in an interior Whitney ball of radius comparable to \(\epsilon\).
Indeed the mixed GFF has the reflected log kernel and the change from Wick normalization at these boundary-scale averages is asymptotic to the same regular-part correction up to a kernel-dependent constant; use standard subcritical GMC uniqueness by mollification (or the reflected Dirichlet field construction and the same subcritical approximation theorem).
Subcritical approximation is in \(L^1\) for the mixed GFF; continuous and normalization corrections give convergence in probability for the fields here.
Lebesgue log-scale averages of these approximants over \(\log(1/\epsilon)\in[N,2N]\) converge in probability too (use \(L^1\) before corrections, transferring by localization/truncation), and integrals divided by \(N\) over bounded-width bands about the endpoints vanish in probability.
All are local assertions within regular open intervals on the boundary.

\proofstage{Transport between collars}

If charts on a common collar differ by a possibly random analytic map of straight boundary segments, with nonzero derivative, localize it on countably many compact subarcs with analytic extension and deterministic smooth bounds as needed there.
Transport a smoothed test kernel to the other chart by change of variables.
It agrees, up to field evaluation error \(o(1)\), uniformly on these localized subarcs, with the standard kernel at radius \(\epsilon\) times the local magnification (using positively oriented flattenings with inward normal up).
To see this, in rescaled coordinates the test density after pullback has the usual kernel as leading term, an \(O(\epsilon)\) linear combination of a bounded family of finitely many smooth derivative-polynomial kernels at the magnified size (Taylor expansion), and a compactly supported smooth remainder \(O(\epsilon^2)\) with fixed arbitrarily many derivatives.
For local mixed Neumann GFFs evaluations of the first-order kernels at centers in a bounded interval and bounded magnification range are \(O(1+\log(1/\epsilon))\) with probability tending to one simultaneously for all small \(\epsilon\): log variance at most \(C(1+\log(1/\epsilon))\), Gaussian tails and unit rescaled parameter boxes with smooth increment bounds, then a union bound on dyadic levels.
The remainder is \(o(1)\) even by local \(H^{-s}\) of the reflected field for small \(s\in(0,1)\), since a rescaled unit-mass smooth density costs \(O(\epsilon^{-1-s})\).
Continuous corrections cause no problem.
The estimates likewise cover random compact subarcs by exhaustion.
Bounded log magnification then disappears in the log-scale averages by positivity and the endpoint-band observation, even when it varies along the line.
The smooth charge shift remains: if the fields are related by \(h\mapsto h\circ f+Q\log|f'|\), the residual factor against the pulled-back measure is \(|f'|^{\eta Q/2-1-\eta^2/4}\); extra continuous fields just give the exponential weight.
This proves agreement for \(Q=Q_\gamma,\eta=\gamma<2\).
\proofstage{Critical passage and coordinate-independent versions}

At criticality use \(\eta\uparrow2,Q=2\) and then the normalized parameter limit in each chart.
The same argument, easier without the straight boundary condition (radially symmetric smooth interior kernels), applies to bulk masses on common conformal neighborhoods.
These comparisons require both fields/masses in the individual charts to satisfy the stated approximation facts; they don't require a theorem of simultaneous covariance for arbitrary laws on both sides of a random coordinate change.
Missing finitely many marked points are handled by diffuseness.
To view a measure as a field function on abstract standard disks, one can integrate versions over Hausdorff charts parameterized by disk automorphisms with a full Haar-class density, keeping the a.e. common pulled-back value; fixed automorphism covariance and Fubini in the standard disk laws give agreement.
Here only the measure \emph{class} in averaging charts matters.
Thus this function may be made exactly coordinate independent on a saturated set of full disk law.
Aperture coordinates on non-Jordan domains, when used, always refer to prime ends.

\end{proof}

\subsection{Subcritical exploration and its marked offspring}\label{proxy:one-branch}

Use in a doubly boundary marked disk the independent symmetric side-swapping chordal \({\rm SLE}^0_\kappa(\kappa-6)\) exploration of a simple \({\rm CLE}_\kappa\), from the first mark to the second.
We use the established Sheffield and Miller--Sheffield--Werner (MSW, CLE percolations and simple CLE on LQG) results on this exploration: loop excursions trace the encountered simple CLE loops once (with left/right signs), the remaining path is the continuous trunk, the full path and trunk can be continued from stopping tips on the trunk with component Markov law (viewed in prime disk coordinates), and the original-target branches couple identically up to separation of targets, with unexplored CLEs in cut-out regions as usual independent copies.
Full nested CLE uses fresh CLEs inside loops as well.
These statements concern \(\kappa\in(8/3,4)\) in simply connected disk/half-plane conformal type, independent of the Liouville field conditional on chart and marked endpoints before running the quantum clock.
We write bSLE for this \emph{branchable SLE} exploration, using the symmetric principal-value convention. Its trunk is the \emph{conformal percolation interface} (CPI); the full branch also includes the completed loop excursions.
The remaining region at a trunk time is the target component \emph{after} any event (on discovering a loop traverse all of it instantaneously for this clock).
Call the two arc lengths from the target back to the post-event tip, keeping consistent side order, \(Z=(Z^1,Z^2)\), \(S=Z^1+Z^2\) \citep{MSWTrunk,MSWQuantum}.

Here are the public LQG exploration inputs used. MSW's quantum half-plane theorem (weight 2 wedge, same simple \(\kappa=\gamma^2\), general skew theorem at skew 0) gives, in quantum trunk time, independent stable boundary change processes in the two sides and four independent Poisson families of \emph{marked} quantum disks: negative jumps cut off disks on the respective side and positive jumps discover loop interiors; perimeter is absolute jump size, boundary root is the pinch (loop attachment or disconnection point) with the once-boundary-marked disk law.
The unexplored surface with tip and target at deterministic trunk time is a fresh quantum half-plane independent of the explored families (MSW half-plane Markov property, Proposition~4.1 in that description).
We only keep the stated pinch mark on offspring.
Choose deterministic units for time in which each base stable process has negative Lévy density \((1/2)|x|^{-1-r}dx\) and positive density \((1/2)(-\cos\pi r)x^{-1-r}dx\); the processes are strictly stable without added location drift.
Kammerer's Proposition~3.6 (version~5) (MSW disk law) gives for the two arc lengths in a \(\gamma\)-quantum disk with length-typical boundary endpoints conditioned on split \(Z_0=(x,y)\), the path law on \([0,t]\) on survival to \(t\), equal to the same stable pair started there restricted to positivity and weighted by \((S_0/S_t)^{1+r}\).
The correspondence of jumps, their sides and geometric cut/discovery is with the same exploration; deterministic choices of length/time units implement the above normalization.
In particular clock can equivalently be recovered up to the known deterministic factor by asymptotically counting discoveries above vanishing length threshold on trunk intervals (stable small-jump counts and absolute continuity).
All uses here are subcritical and in exactly disk or weight 2 half-plane law.
We will include pinch marking and post-disk factorization for disks; we explain the transfer below instead of requiring any critical marked-cut law \citep[Theorems~1.4--1.5, Section~3.1, Proposition~4.1, Corollary~4.2, and Proposition~4.4]{MSWQuantum}\citep[Proposition~3.6, version~5]{Kammerer}.

\begin{proposition}[Marked disk factorization]\label{proxy:marked-factorization}
For almost every prescribed positive initial perimeter pair, run the independent symmetric $\SLE^0_\kappa(\kappa-6)$ branch on the corresponding boundary-marked disk.  At a fixed surviving trunk time $t$, conditional on the signed boundary-length path up to $t$, the offspring discovered by $t$ are independent once-boundary-marked disks with perimeters equal to their absolute jump sizes.  Their mark is the actual attachment or disconnection point.  Independently, the remaining surface with its tip and target is a twice-marked disk conditioned on its two arc lengths.  The assertion concerns abstract marked surfaces and retains the component laws of the unobserved CLEs.
\end{proposition}\begin{proof}

\proofstage{Strip likelihood}

Randomize horizontal translation in the strip with Lebesgue measure.
Consider first \(\sigma_W\), radial part \(m\) with two-sided variance-2 Brownian increments with drift \(+a\), \(m(0)\) Lebesgue independently; and \(\sigma_D\), changing this to initial intensity \(e^{-a m(0)}dm(0)\), independent variance-2 Brownian increments \emph{away} from 0 on both sides with drift \(-a\).
Both get the independent lateral field.
These are translation invariant, and on the strip up through real coordinate \(v\) their likelihood ratio is
\(\exp(-a m(v))\).
Modulo translation they are, up to constants, the DMS quantum half-plane measure (a multiple of probability) and the twice-marked disk measure.
For example divide by translations using a unit-location cutoff on first reach of average height \(k\).
For \(\sigma_W\) this has finite positive mass (integrate the initial Lebesgue height against Brownian first crossing, or use the bound by the expected one-unit range and transience on the left), independent of \(k\).
The path after first reach of \(k\) has the unconditioned forward increment law by the strong Markov property (localize to finite windows); for the past center first at first reach of \(-K\), use translation invariance of the quotient, and let \(K\to\infty\) in the ascending Brownian segment, reversing at \(k\) as in \eqref{eq:original-1} to obtain its deficit \(U_a\) before first reach.
This is the usual first-exit embedding of the weight-2 wedge (\(\gamma\)-insertion parameter).
For \(\sigma_D\) restricted to reaching \(k\), transfer the path to this stopping time by the likelihood \(\exp(-ak)\) (on bounded time windows by Brownian optional stopping), then use the forward drift \(-a\).
This is exactly the description of the scaled log Bessel excursion up through a fixed level (upward part \({\rm BES}(4-\delta)\) ascent, followed by unconditioned \({\rm BES}(\delta)\) descent to zero in Bessel units, \(\delta=3-4/\gamma^2\)).
Vary \(k\); independent lateral projections remain stationary.
Thus after dividing by translations the \(\sigma_D\) perimeter-pair intensity is a constant times \(S^{-1-r}dx\,dy\) by the maximum and marking rules above \citep{DMS}.

On \(\{T>t\}\) where \(T\) is duration of the branch, map the strip to the remaining region, mapping ends to tip and target by \(F\) with \(F(u)-u\to0\) at the right end (which has an untouched collar).
This difference and all its derivatives in the tail are \(O(e^{-\Re u})\) by reflection in exponential coordinates.
Denote the charge-transformed field by \(\tilde h\), its average by \(\tilde m\).
Under \(\sigma_W\), MSW gives the product of the half-plane cut Poisson law up to \(t\) (including the signed length changes) and \(\sigma_W\) for \(\tilde h\), with \emph{exact} constant one: at the quotient level this is the MSW normalized probability assertion, and the gauge of \(\tilde h\) is equivariant under parent strip translation, hence has the same translation factor \citep[Proposition~4.1, Corollary~4.2, and Proposition~4.4]{MSWQuantum}.

\proofstage{Domination}

We claim domination for \(\sigma_D\) on \(T>t\) of these same outputs, by the cut Poisson law times \(\sigma_D\) for \(\tilde h\).
Test a nonnegative functional of all cut marked disks and signed changes up to \(t\) and \(\tilde h\) up to horizontal level \(K\).
Localize temporarily to the event that the hull up to \(t\) and these observed field restrictions come from before a large horizontal level \(v\).
This can be tested there, with independent full bSLE: run up to leaving a slightly smaller tail cutoff, using counts of discovered loops and their quantum lengths to read the clock, requiring that trunk time already exceeds \(t\).
Any partial cutoff excursion can be discarded here.
Indeed the count convention recovers the same duration under both measures on all completed portions by the public laws above, simultaneously by testing rational chronological bounds.
Loop lengths on portions compactly before the cutoff are read on their insides using the field restrictions and canonical univalent charts; under the wedge they are disk lengths, and the same locality holds under the disk by finite-cutoff absolute continuity (or the MSW unmarked disk rule).
The signed change paths themselves can be read from the jumps by compensated stable summation on positive-survival compact intervals, using the same subcritical path laws.
Thus there are such field-before-\(v\) tests exhausting the survival event as \(v\to\infty\).
In this use of a measure likelihood all event/test definitions need only agree within their local measure classes.

With \(\chi_v\) a smooth probability kernel on \([v-2/v,v-1/v]\) require also that \(F\) maps this slab before \(v\) and
\[
 |m(v)-\int\tilde m(y)\chi_v(y)dy|\le\epsilon .
\]
Along integer \(v\) this extra localization is eventually immaterial under \(\sigma_D\) on survival.
Indeed local fixed-negative-order distribution norms of the lateral field on unit slabs, including the Neumann edges by reflection, are eventually polynomial, as are radial suprema; \(F\) is exponentially close to identity in every smooth norm, so changing the averaged test density on the slab (of only polynomial smooth norms) and adding the charge correction causes a vanishing error.
Brownian oscillation on the \(2/v\)-window vanishes along these integers.
By Fatou and the parent likelihood the \(\sigma_D\)-integral is bounded by \(\liminf\) of \(\sigma_W\)-integrals with the likelihood bounded on the localization by \(e^{a\epsilon}\int e^{-a\tilde m(y)}\chi_v(y)dy\).
Drop localization on the latter side; Fubini and the precise wedge product give \(e^{a\epsilon}\) times the claimed product integral by the same Brownian likelihood in the remaining strip (where field tests use only up to \(K\)).
Send \(\epsilon\) to zero.
Cylinder domination extends to the full output by exhaustion/monotone class (use finite product-mass cylinder cutoffs, e.g. restricting the remaining average and any needed translation gauge).

\proofstage{Disintegration}

Divide by translations using a unit cutoff on the remaining disk gauge.
In the dominating product the putative \(Z_t\) is the remaining split length; in the survival law that equality and the signed change law hold by the definition and the cited disk theorem.
Keep only paths with valid positive backtracked pair \(Z_s\), \(0\le s\le t\).
On projecting just to the absolute lengths/path, this restricted dominating mass is exactly saturated by Kammerer's law: indeed its density with respect to stable changes and \(dx\,dy\) at time zero is the disk constant times \(S_t^{-1-r}\) by translating the two lengths.
The disk constant is the same one as for the initial disk.
Thus the domination is equality.
Disintegration proves the full marked-factorization rule for almost all initial split lengths: given the path up to \(t\) on survival, the offspring are independent once-marked disks at their perimeters, and independently the remaining surface with tip and target is a twice-marked disk conditioned on the two arc lengths.
This assertion concerns abstract marked surfaces, not independence of embedded maps.
All such explorations may carry along the unobserved CLEs with their component Markov law, by first conditioning on the branch/CPI geometry and field \citep[Proposition~3.6, version~5]{Kammerer}.

\proofstage{Clock and stopping-time conventions}

A few details about times and measurability in this transfer.
All laws of surfaces/cut disks in the MSW input are modulo conformal re-embedding, retaining the indicated pinch points as boundary marks; a possible additional mark/record of original boundary on a negative disk is not needed as another independent sample.
Clock constants in wedge and disk can separately be calibrated exactly by the number of completed positive excursions whose \emph{inside boundary} has length in shrinking bins (in the same deterministic length units).
In particular their calibration agrees on any common portion under the local likelihood comparison, without referring to how much absolute length there may be far to the right.
Indeed the stable-process law of large numbers for those bins (e.g. almost surely along a fast subsequence and at rational trunk times, then by monotonicity) applies also up to each fixed compact surviving positive path horizon by the disk path law.
This gives simultaneous elapsed trunk time on all observed portions, with the clock flat during any loop tracing in expanded time.
Absolute jump marks can likewise first be tested on countably many lists at positive thresholds, and the stable change process recovered by compensation (including the strictly stable centering).
Boundary measures on already completed disk interiors are local field functions in those disk charts.
Thus in using the parent likelihood one may require common observed clock \emph{strictly beyond} \(t\) before the cutoff and rational slightly larger hull bounds, discard incomplete excursions there, and exhaust; definitions of the outputs up to \(t\) agree on these tests under the two parent laws.
For any surviving bounded initial piece all needed earlier trace and completed interiors fit before a sufficiently far cutoff.
Only the full independent exploration geometry, not field data later than the cutoff, is allowed for free in this likelihood.
For the quotient interpretation of \(\sigma_W\), one could equally change to the gauge of the \emph{last} reach of 0, by translation and reversal with sign (symmetries of its radial model); then the free and barrier-conditioned halves occur in the opposite order.
These are the two exit descriptions of the same drifted wedge, and only the translation quotient matters \citep{MSWQuantum}.

Continuation of the \emph{full} exploration at the zero-driver-gap trunk times considered here (including just after a disconnection or completion of a selected loop excursion) uses the target-component bSLE Markov property.
One can regard all parent field/end data as known initially in applying geometric strong Markov in the independent exploration; tip times defined by the length process on observed paths are adapted using the count clock or its right-continuous inverses as above.
For example first jump times past fixed thresholds can be tested by disconnections and \emph{completed} excursions with their component marks/lengths, and deterministic trunk horizons can be located by when completed-path trunk time first exceeds the prescribed value, in right-continuous filtration.
We never ask for a fresh full bSLE at a stopping tip in the \emph{middle} of tracing a loop.
Future unrevealed CLE geometry in the components has the corresponding CLE-exploration conditional law.
These uses are the ordinary domain Markov/branching and zero-gap restart properties of the MSW/Sheffield bSLE construction, combined with the independent field, in disk or half-plane conformal type as above \citep{MSWTrunk,MSWQuantum}.

\end{proof}

\begin{lemma}[Jump times and prescribed limiting starts]\label{proxy:jump-times}
The factorization of Proposition~\ref{proxy:marked-factorization} remains valid just after length-path stopping jumps, on compact positive-survival localizations.  Every prescribed positive critical starting pair can be approximated by good subcritical starting pairs, using standard disk charts based on the source and an additional independent length sample.
\end{lemma}\begin{proof}

\smallskip\noindent\textit{Stopping just after jumps.}\quad
One may stop just after jumps, selected as stopping times of the length path (including crossing specified absolute jump thresholds).
Approximate by dyadic times from above on survival.
In the canonical prime disk of the true remaining region at that jump, the subsequent hulls shrink to the post tip.
Indeed clocks parametrize the continuous trunk, attached loop excursions are locally finite at positive diameter, and continuation in this disk is continuous at its starting prime end by bSLE Markov at the jump (this is a stopping time for the exploration given the independent parent field, or use the completed loop-return time for a positive jump).
Normalized maps to the slightly later domains \emph{in this disk} (normalizing at 0 or another fixed interior point) converge to identity uniformly up to the boundary: reflected convergence on boundary arcs off the tip, and on the remaining sector confinement by small crosscuts about that tip, or the maximum principle.
Pullback charts of the later remaining marked disks, using the post tips there and added independent boundary marks for standard twice-typical coordinates, differ from these maps by automorphisms which cannot escape compacta in probability on compact ranges of the two lengths.
Indeed an open arc about the original target is untouched, with the same length by the arc comparison in later marked charts.
Choose a third comparison point in that arc a sufficiently small fixed length away from the target (exhaust by countably many such lengths).
Its locations, the targets and tips stay separated in the true remaining canonical disk, while their preimages in the standard twice-typical charts are separated in probability by diffuseness/tightness of the boundary measures and by their positive length spacings. A degenerating automorphism would collapse two.
Here at the limiting jump the positive split values are supplied by the signed càdlàg paths; their identification with actual post lengths is not presupposed.
Interior distribution pullbacks converge through the normalized maps since all derive from the same parent field.
Thus by subsequences and the dyadic laws the true post field and marks in prime coordinates have the same asserted law at the limiting split, including true post lengths; actual post tip and target are the prime limits through these shrinking continued hulls.
One obtains joint independence with all cuts up to that jump by the same test limit at dyadic times (disk laws with prescribed splits vary continuously by common twice-typical coordinates and length scaling/labels; keep first finitely many already completed jumps and exhaust).
Stopping tests can be localized to compact positive survival; boundary or axis hits not under this assertion.

\smallskip\noindent\textit{Prescribed starting pairs.}\quad
For the single-branch subcritical law we only use almost everywhere initial perimeter/split prescriptions (or AC randomizations of them).
In deriving the law later for any \emph{critical} prescribed positive length pair, one simply approximates the prescribed start by such good subcritical pairs, or by compactly supported diffuse laws on the pairs tending to the prescription.
The parent is a once-typical-mark disk at this perimeter with the target at the relative oriented length label given by the pair; for bounds or taking limits embed using an extra fresh independent typical boundary point with the source, \emph{not} using the target itself as the extra point.
Thus given the start lengths this chart has the perimeter-scaled twice-typical field and genuinely has the convergence and estimates above/below uniformly for converging positive pairs, including convergence of the target by diffuse full-support chaos.
Conditional on that data the unclocked bSLE/CLE is just the independent conformal-domain law with the indicated endpoints.
Almost-sure properties requiring, for example, other fields/charts simultaneously can first be proved by averaging at each subcritical parameter before choosing the approximate starts.
No uniform estimate in the widths of small split intervals by dividing general error probabilities by those widths is needed.

\end{proof}

\subsection{Strict boundary-label inheritance}\label{proxy:strict-stack}

Figure~\ref{fig:stack-labels} illustrates the strict label match used below.
\begin{figure}[htbp]
\centering
\begin{tikzpicture}[x=1.12cm,y=1.02cm,>=Stealth,
 every node/.style={font=\small}]
\draw[->,black!65] (0,0) -- (6.05,0) node[right] {trunk time};
\draw[->,black!65] (0,0) -- (0,3.55)
 node[above,align=center] {target-backward\\label on side $i$};
\node[below left] at (0,0) {$0$};
\fill[blue!6] (1.5,1.7) -- (1.5,2.7) -- (2.15,2.43)
 -- (2.8,2.57) -- (3.45,2.25) -- (4.35,2.35) -- (4.35,1.7) -- cycle;
\draw[very thick,black!70] (0.22,1.3) -- (1.5,1.3);
\draw[very thick,orange!85!black,->] (1.5,1.3) -- (1.5,2.7);
\draw[very thick,black!70] (1.5,2.7) -- (2.15,2.43)
 -- (2.8,2.57) -- (3.45,2.25) -- (4.35,2.35);
\draw[very thick,red!70!black,->] (4.35,2.35) -- (4.35,0.8);
\draw[very thick,black!70] (4.35,0.8) -- (5.65,1.02);
\draw[densely dashed,black!40] (0,1.7) -- (1.5,1.7);
\draw[thick,blue!75!black] (1.5,1.7) -- (4.35,1.7);
\fill[blue!75!black] (1.5,1.7) circle (2pt);
\fill[blue!75!black] (4.35,1.7) circle (2pt);
\node[left] at (0,1.7) {$h$};
\node[below left,font=\footnotesize] at (1.47,1.3) {$Z^i_{s-}$};
\node[above left,font=\footnotesize] at (1.5,2.7) {$Z^i_s$};
\node[above right,font=\footnotesize] at (4.35,2.35) {$Z^i_{t-}$};
\node[below right,font=\footnotesize] at (4.35,0.8) {$Z^i_t$};
\node[below,blue!65!black,font=\footnotesize] at (2.925,1.65) {horizontal match};
\node[align=center,above] at (3.3,3.15)
 {$Z^i>h$ between the apertures};
\node[orange!85!black,align=center,below] at (1.5,0.95)
 {inserted loop\\boundary};
\node[red!70!black,align=center,right] at (4.55,1.8)
 {detached\\disk boundary};
\node[below] at (1.5,0) {$s$};
\node[below] at (4.35,0) {$t$};
\draw[black!40,dotted] (1.5,0) -- (1.5,1.25);
\draw[black!40,dotted] (4.35,0) -- (4.35,0.75);
\node[align=center,font=\footnotesize] at (2.9,-0.95)
 {Arrows show jump direction; boundary coordinates increase upward.\\
 Each aperture's two endpoint labels map to its own root pinch.};
\end{tikzpicture}

\caption{A finite schematic of the boundary-label rule. At time $s$ a positive jump inserts a loop boundary; at time $t$ a negative jump removes a disk boundary. The level $h$ lies in both open apertures and strictly below the intervening extended graph, so the two marked labels represent the same boundary point by Lemma~\ref{proxy:stack-rule}. The endpoints of each aperture map to that aperture's own root pinch; no identification of the two apertures' endpoint pairs is asserted. Arrows show jump direction, while boundary coordinates increase from the lower endpoint of either aperture. The drawing represents neither a typical Cauchy trajectory nor Euclidean lengths or shapes of the welded surface.}
\label{fig:stack-labels}
\end{figure}

\begin{lemma}[Strict stack rule]\label{proxy:stack-rule}

For the almost-everywhere subcritical laws of Proposition~\ref{proxy:marked-factorization}, the following boundary-label inheritance rule holds.
Denote the remaining disks by \(D_t\), jump disks by \(C_t\) and the trunk tip by \(P_t\).
Use canonical prime coordinates in a starting post-state, at an ordinary deterministic trunk time (including 0) or at one of the above jump stopping times.
Boundary maps of remaining regions and offspring extend continuously in the original closed disk (the traces and their initial segments are continuous curves, have locally connected unions with the boundary, and the components have locally connected boundary).
In the starting disk's own coordinates subsequent bSLE branches are the ordinary continuous explorations, with continuous closure maps there as well.
Consider one side's boundary prefix from target back toward tip to length label \(h>0\).

\begin{enumerate}

\item Through a time interval on which the extension of that side's graph of lengths (including vertical jump segments, except possibly at specified endpoints) stays strictly above \(h\), the prefix is unhit, in correct \emph{access} in the subsequent prime frontiers.
On shared untouched open arcs its parametrization is inherited.

\item A positive jump inserts the full loop boundary interval from the pre- to the post-value in target-backward labels, with lengths as on the corresponding jump disk, cyclically from its root in the prescribed sense (so both endpoints have image the pinch).
At a negative jump, the interval from post- up to pre-value parametrizes in the corresponding sense the detached disk boundary starting and ending at its root.
Here the negative jump aperture notation uses length \emph{on that detached disk}; a comparison to earlier post-frontiers is on transmitted portions (it does not require a general continuity-of-boundary-chaos theorem on a pre-jump kernel limit).
If a strictly transmitted prefix has labels falling in the \emph{open} negative jump interval at its first loss, those labels are on this detached boundary with coordinate given by subtracting the post-value.

\end{enumerate}

Consequently strictly horizontal matches (staying strictly below the intervening path, between interior aperture/prefix labels) refer to the same actual points, with the stated prime-side inheritances, not just matching total masses.

\end{lemma}\begin{proof}

\proofstage{Protected prefixes and negative losses}

Work just after the specified start in its own canonical disk.
The target has an untouched neighborhood on every pre-terminal compact path segment.
Before any first intrusion, untouched open arcs use exactly the length inherited from this disk by collar comparison. A first hit into the interior of such an arc, coming from outside, separates the rest toward the tip on that side, making a disconnection event, i.e. a negative jump.
To formalize for a protected prefix, test with an intermediate endpoint \(h'\) slightly \emph{farther} than \(h\), still strictly below the graph; such test points can be chosen untouched.
Indeed in starting prime coordinates the continued exploration has the ordinary bSLE law, for which specified other boundary points are not hit (only the trunk could hit, as complete future loops are simple CLE loops).
Here one can use MSW's trunk law at skew zero, chordal \({\rm SLE}_{16/\kappa}(8/\kappa-3;8/\kappa-3)\), and the standard SLE with force points boundary-intersection theorem (in particular the boundary dimension bounds of Miller--Wu and related works): away from start and target, for an initial same-side weight strictly above \(\kappa'/2-4\), with \(\kappa'=16/\kappa\in(4,6)\) and the two tip weights \(>-2\) here, it is not boundary filling and has zero boundary Lebesgue measure there \citep[Theorem~1.6]{MillerWuBoundary}.  Indeed the boundary dimension at $\rho=\kappa^\prime/2-3$ is $1/2+1/\kappa^\prime<1$.
The one-side local bounds suffice (near a boundary swallow on that side localize away from opposite force points); Möbius covariance then gives avoidance of a given other point.
Even empty interior of the hit set suffices to choose untouched intervening labels, since boundary chaos in prime coordinates has full support without atoms.
These continuations need only be used from the countable post-state times under consideration. A first hit on the arc from target to \(h'\) then cuts off a nonempty boundary interval including \(h'\) on that side.
The post tip in the target component is at this hit with the access on the target side; that side length is exactly the coordinate of the hit.
For the last assertion, along the arc from the target the post domain has untouched access all the way up to the hit, converging there to the tip prime end, and in the cut domain the just lost part has access from the other direction also tending to its pinch prime end.
There is no extra split into different access labels at the hit in either component: in these starting-disk coordinates the prior trace avoids the hit until this first contact, is connected to the earlier start, and any two inequivalent accesses there joined by a simple interior curve from and to the point would bound a pocket in the disk with no prior trace inside (the trace is disjoint from that joining crosscut and the point until its terminal visit, and starts outside).
Hence the pocket lies in the component; equivalently the boundary interval between those accesses on the pocket side in Riemann coordinates would have constant image, impossible for a univalent map.
This also follows by trimming away from the approaching final trace and the untouched side arc.
Full bSLE excursions before contact are included in the connected trace argument.
Positive loops never intrude onto the starting boundary when they are in a fresh starting disk.
If a drop crosses \(h\) while a slightly longer prefix was preserved up to just before it, it must hit into that longer arc at the drop: otherwise the prefix would still give the lower bound (a disconnection swallowing the arc also requires such a hit between it and target).
The disconnection point's post length gives the claimed difference coordinate in the negative disk by the same shared arc rule.
This argument uses disconnection-post and cut markings at the jumps, covered by countable size-threshold lists \citep{MSWTrunk}.

\proofstage{The inserted loop and its two boundary measures}

For a positive discovery, the loop is attached at the trunk and in the post boundary, walking back from the tip on the insertion side traverses just the inserted loop first (two copies of the pinch); the interior is the corresponding \(C_t\).
These are the excursion and side conventions of bSLE, the loop other than its attachment being away from the previous compact trace.
More explicitly complete positive-excursion loops discovered strictly later than an ordinary surviving trunk time are still simple CLE loops in the unexplored target domain at that time, by the bSLE exploration property in those coordinates.
Use also countably many such times increasing toward the discovery; earlier contacts even at an eventual pinch would be excluded on these earlier pieces.
The only accumulation from the immediate approach is at the attachment by path continuity.
There is no simultaneous other detached open component at that time by the cut/jump correspondence and the Poisson/stable laws (two such positive-perimeter events do not coincide).
Thus away from the pinch along the simple loop the exterior access belongs to the post frontier, and the two disk sides there have local collars for comparing this actual seam \citep{MSWTrunk,MSWQuantum}.

Its length agrees on both sides not merely as a total but on open subarcs: use the DMS/Sheffield subcritical chordal quantum zipper \citep{DMS}\citep[Theorem~1.8 and Remark~5.7]{SheffieldZipper} (ordinary independent chordal \({\rm SLE}_\kappa\), \(\gamma=\sqrt\kappa\), weight-4 wedge into weight-2 wedges), which identifies lengths on the two sides.
This transfers locally to the present loop away from its start/end.
Indeed, viewed conditionally from rational capacity times during its tracing (also before any prescribed strictly interior subarc), bSLE until the excursion closes and away from the return point is ordinary simple chordal SLE with a boundary force/drift, locally absolutely continuous up to localized stops with separated driver points.
Its fresh open segment is interior in the then domain.
The parent field is independently, on interior compact subsets, a GFF plus a continuous correction up to change of law, and can there be compared by Markov coupling to the weight-4 zipper input (difference continuous on the neighborhood; local wedge chart and both fields couple via the same Dirichlet projection).
Condition first on the geometry up to the initial time for this comparison, then localize to subarcs in those neighborhoods and stopping bounds for ordinary chordal absolute continuity.
Zipper equality there concerns only that open part of the seam: completed ordinary-SLE collars from each side carry mixed-Neumann plus continuous regular arc charts, and other completions with the same local germ, including the actual loop cut/post charts, use exactly the same measures up to the exponential weight of the continuous field correction on the seam, by arc comparison.
One can phrase the localized transfer just as an assertion about smooth-kernel approximations on common side collars and thus needs no absolute continuity for the rest of the loop.
To detail that local comparison, condition at a geometry-only initial capacity time in an excursion, use an interior domain library for compact segments lying ahead away from the force coincidence, and couple the parent and zipper fields on the chosen neighborhoods by the same interior zero-Dirichlet projection there, first in reference field laws restoring independent average processes.
Those projections can be taken identical; the discarded harmonic fields, average substitutions and conformal terms are continuous in each interior neighborhood, and the disk bias can be included by absolute continuity.
Thus the zipper copy still satisfies its length equality there a.s.
One can first draw a common ordinary SLE segment independently, transfer up to localized separated-driver stops by chordal SLE drift absolute continuity, and then complete the actual exploration (the zipper conclusion was local to the compared seam with its own completion).
Simple-arc subcollars share the same side as these completions and use only those field restrictions, with the common continuous-surface correction on the seam.
Actual post and cut arc charts qualify for the earlier arc comparison by their disk marginals just proved; zipper charts away from their end marks similarly by wedge laws.
Countably many such compact-segment localizations cover the open loop arcs.
This applies first in typical-target coordinates, where the whole branch is independent of the parent disk field before clocking, to every excursion indexed using also field-independent rational times; almost-everywhere statements suffice for the above choice of approximations.
Endpoints have no atoms.
Together with the full signed jump rule this identifies the inserted interval as claimed.
\proofstage{Position and order of the inserted interval}

The open loop excluding the attachment is inside the pre-loop target component with closure based at a fresh tip, and is adjacent externally to the post component as above.
The attachment, fresh even through the immediate approach, has only one access in the pre-loop target component.
Indeed two accesses joined by an interior Jordan arc from and to that point would enclose a pocket missing the earlier connected trace (which starts on the outer boundary and avoids the point until then), again impossible by nonconstancy on prime intervals.
Thus in pre-loop prime coordinates the circle can only add a boundary detour from and back to that tip, with the post tip at one of its ends.
Its open arc carries equal interior/exterior length by the preceding comparison, so the detour at the top of that post side has exactly the jump magnitude.
The rest of each post side, before this detour as applicable, has length at least the corresponding pre-value: indeed work in a sufficiently close preceding post-state (ordinary times increasing to the event suffice), where prefixes strictly short of those pre-values stay protected up to insertion and persist still outside the detour since the new loop is interior relative to that state.
Total length increase now gives equality and the side/order convention as claimed.
Horizontal comparisons from a positive aperture just start in its post-state with inherited labels.

\end{proof}

\subsection{A strictly subcritical metric proxy}\label{proxy:metric}

To control actual access maps (mere interior kernel convergence is not a priori enough), fix an auxiliary strictly subcritical LQG metric with parameter \(\gamma_0=\sqrt{8/3}\), charge \(Q_0\), metric exponent \(\xi_0=\gamma_0/4\).
We use the subcritical metric construction and its LFPP realization \citep{GMMetric}, together with the weak-metric moment estimates of Dub\'edat--Falconet--Gwynne--Pfeffer--Sun \citep[Theorem~1.9 and Proposition~3.1]{DFGPS}.  The inputs are $\xi_0Q_0=5/6<1$, scaling and Weyl scaling, and upper moments of all orders for internal annular/path-around costs and inverse moments of all orders for across-annulus costs on fixed interior annuli in a whole-plane GFF with fixed circle constant.
Fixed internal box-crossing upper estimates used below follow by chaining finitely many such short winding paths over overlapping interior annuli with forced intersections.
The theorem of Devlin on simultaneous conformal coordinate change is used \emph{only for this strictly subcritical proxy}.
In a fixed bounded parent disk coordinate \(\Delta\) (unit disk), use \(D^{(0)}_{h_{\rm parent}}\) as local/intrinsic proxy metric.
If \(F:{\cal S}\to U\subset\Delta\) is a jump disk or jump-post-state map in standard marked coordinates as above, the internal proxy pullback has length element that of \(D^{(0)}_{g}\) weighted by \(|F'|^{\xi_0(Q_0-Q_\gamma)}\), where \(g=h_{\rm parent}\circ F+Q_\gamma\log|F'|\) is the actual disk field.
Since \(Q_0>Q_\gamma\) here, Cauchy estimates give the deterministic upper weight bound \(C(\min(1,\operatorname{dist}(\cdot,\partial{\cal S})))^{-\xi_0(Q_0-Q_\gamma)}\).
The same applies to compositions of such maps into the fixed parent chart \citep[Theorem~1.9 and Proposition~3.1]{DFGPS}\citep[Theorem~1.1]{Devlin}.

Here are details about the randomness in this metric comparison.
For a branch toward a typical target choose parent charts from its start and target, with deterministic strip-to-disk convention.
Conditional on the full \emph{unclocked} bSLE branch, independent of the parent field, all countably many domains of jumps and jump-post states can be treated by ordinary fixed-map covariance, mapping first from a canonical disk by geometry-only choices.
Indeed they come from the geometry of loop closures and disconnections alone.
Conditional on it, on each domain the parent field and pullbacks qualify in law up to absolute continuity as a whole-plane GFF plus a (random) continuous function in its interior: use the lateral strip description, add back and remove an independent average, use Neumann/Dirichlet Markov comparison (free GFF has independent Dirichlet projection plus harmonic distribution in the \emph{open} subdomain), then whole-plane/Dirichlet comparison and charge correction; normalization/tilt is continuous addition/change of law.
After fixed geometry-only comparison, the remaining h-dependent re-embeddings change only standard disk coordinates by conformal maps between \emph{fixed} domains, for which use Devlin's simultaneous rule on the disk.
At the random standard chart on the child side the field with proxy charge is \(g+(Q_0-Q_\gamma)\log|F'|\), with the same qualification now by the marked disk marginal law; its metric and that for \(g\) therefore are related by the usual continuous Weyl length factor even here.
To clarify this last version argument, the subcritical metric on each fixed domain in these laws agrees locally with the normalized LFPP limit length element (one can use the localized dyadic approximants).
For a whole-plane GFF plus continuous this is the subcritical LFPP convergence/metric identification, and on any coupling having two such qualifying marginal laws differing by a random continuous function the two elements of length differ by the exponential weight: locally the approximating mollifier weights sandwich with arbitrarily small oscillation error.
Locally uniform distance convergence suffices for this comparison, by first confining short distances within compact neighborhoods via the positive exit cost, passing the sandwich for nearby pairs, and subdividing compact paths.
One may equivalently apply the pathwise continuous-addition construction of the metric from LFPP; convergence/identification on local compact neighborhoods transfers under absolute continuity.
This uses Devlin as a \emph{version of} the metric on the fixed standard domains, avoiding evaluation of a version on a field-dependent domain outside its scope.
Thus arbitrary changes between qualifying standard parent charts for these branches cause only the corresponding analytic/Weyl terms.
In particular the a.s. comparisons transfer also if using a fresh typical mark with the source for the parent chart, and then for almost all the above subcritical perimeter/split prescriptions in this chart.
In successive cuttings apply it one step at a time in the marginal parent/branch laws, transporting earlier derivative weights by path composition.
We will use this also directly at critical charge 2 only for analogous branch laws when the independent geometric branch and actual marked offspring laws there have been established (the proxy still subcritical) \citep[Theorem~1.1]{Devlin}.

\begin{lemma}[Uniform lower cost for macroscopic paths]\label{proxy:lower-cost}

In a unit disk parent chart obtained from standard strip coordinates with typical source and an auxiliary or target typical mark and perimeter 1, paths in the open disk of macroscopic Euclidean diameter cannot have arbitrarily small internal proxy length, uniformly in probability as \(\gamma\) increases.
The assertion is uniform when the perimeter belongs to a compact subset of $(0,\infty)$. 

\end{lemma}\begin{proof}

On interior compacta this is ordinary metric positivity and tight continuous corrections; in bounded horizontal slabs up to the strip edge it suffices to rule out arbitrarily cheap paths close to the boundary crossing a fixed interval tangentially.
For each fixed horizontal coordinate \(x\) on that interval, consider interior Whitney boxes above \(x\) at small depths \(e^{-n}\), with finitely many per level, arranged so any path hitting the normal above \(x\) at such depths and exiting its Whitney neighborhood (size at most \(C\) times depth) must cross one interior annulus of size \(\asymp e^{-n}\).
In the straight local chart their proxy charge scaling is \(e^{\xi_0(H_n-Q_0 n)}\) up to tight coarse continuous factors, times annular costs with uniform inverse moments.
The latter fact follows by rescaling the Neumann field minus semicircle average (bounded harmonic Gaussian corrections on the rescaled interior boxes) and the stated interior metric estimates.
The deterministic parent chart derivative correction on the slab is bounded.
Since \(H_n/n\to0\) in probability uniformly over all large \(n\) at fixed \(x\), with summable Gaussian error tails, and \(\xi_0Q_0<1\), with uniformly high probability at fixed \(x\) all these small crossings cost at least \(e^{-n}\) (the DMS tilt and random coarse corrections can first be truncated).
Fubini gives a good set covering most of the fixed interval in Lebesgue length.
Hits by a path confined to the thin collar and traversing tangentially thus force total length bounded below: select disjoint horizontal Whitney neighborhoods of hits on good normals with sum of widths bounded below by interval covering (use a strictly interior subinterval to ensure exits).
Paths that do not remain in a thin collar already have an interior macroscopic traversal.
Any macroscopic movement on the compact disk boundary contains tests away from the two ends, covered by finitely many such slab tests.
This proves the estimate.

\end{proof}

\begin{proposition}[All-point boundary-chain estimate]\label{proxy:boundary-chain}

For a disk \(g\) in standard twice-typical strip coordinates with perimeter \(s\), paths in its interior approaching any flat boundary point \(z\) and \(+\infty\) can be joined, along the direction to the right, at cost bounded (including the above deterministic upper weight) by

\begin{equation}
K\,s^{2\xi_0/\gamma}\left(\frac{m([z,+\infty))}{s}\right)^{\theta 2\xi_0/\gamma},
 \label{eq:original-11}
\end{equation}
where $m$ denotes boundary length on that side.  The moment quantifier is as follows: for each $M\in(0,\infty)$, choose $\theta<1$ sufficiently close to $1$ and then $\gamma_M<2$ so that
\[
 \sup_{\gamma\in(\gamma_M,2]}\E[K^M]<\infty.
\]
The value $\gamma=2$ denotes the limiting disk law.  The choice of $\theta$ is allowed to depend on $M$.
More precisely we assert existence of connected chains/paths \emph{between arbitrarily small neighborhoods} of the two points, simultaneously at every \(z\).
The reverse statement toward \(-\infty\) also holds, and the chains may approach an end on either side using central Whitney slab paths.
No exact internal metric extension to boundary is needed.

\end{proposition}\begin{proof}

We prove \eqref{eq:original-11}.
On a unit slab use first the mixed/free GFF boundary model above, writing \(H_n(x),X_n=2n-H_n(x)\) with the previous reference.
Cover a boundary interval by mesh points at each scale \(e^{-n}\).
Chains to larger scales above any point of the interval can be made by short crossing and surrounding paths in interior Whitney boxes: for instance place annuli about an interior point at height \(c e^{-n}\) above each mesh center, with radii small multiples of that height, and connect by chains of overlapping annuli in the comparable-depth neighborhood to the possible next coarser centers, using forced crossings and staying at comparable height.
Windings at successive levels and these bridges join; finitely many configurations at each center suffice.
Centers of boxes need not lie directly to the right of the desired point.
Charge the estimate at scale \(e^{-n}\) against the mass of a short boundary interval at displacement \(+C e^{-n}\) on the same side from its grid center, so all charged intervals are within a bounded distance to the right of the limiting \(z\).
Under the endpoint tilt by \(H_n\), the mass divided by \(e^{(\gamma/2)(H_n-Q_\gamma n)}\) has inverse moments of all orders as proved above, and the crossing costs (summing the finitely many tests at this center) with the derivative upper weight, divided by \(e^{\xi_0(H_n-Q_\gamma n)}\), have positive moments of all orders.
Indeed without the derivative weight use the standard rescaled interior metric estimate with \(Q_0\); interior harmonic-field corrections after subtracting the coarse semicircle, including the tilt shift, have uniform exponential moments.
Thus the cost bound at the center divided by the indicated short mass to power \(\theta 2\xi_0/\gamma\) is at most
$Y_n e^{-dX_n}$, with $d=(1-\theta)\xi_0>0$ and $Y_n$ having uniform moments of all orders under the tilt.  Indeed, the exact remaining exponential factor is $e^{-d[X_n+(Q_\gamma-2)n]}$, which is at most $e^{-dX_n}$ because $Q_\gamma\geq2$.

\smallskip\noindent\textit{Uniform control of all chains and visit counts.}\quad

Here are details for summing these charges on all boundary chains.
Use the simultaneous barrier with parameter \(y\ge1\) already proved.
Except for probability \(C_d(1+y)^C e^{-y}\), at all grid points in the band \(X_n\in[-y+j,-y+j+1]\) with barrier survival, \(j\ge0\), we have \(Y_n\le C_d e^{d j/4}\) and oscillation \(O_n\le C+j/3\) up to fixed mesh multiples.
Indeed sum the undo-tilt bounds with cost \(C e^{-y+j}\), using the high-moment/Gaussian tails and saving ballot power \(b>2/3\).
Thus along chains using nearest mesh centers at each scale, ratios are bounded by \(C_d e^{dy}e^{-d(X_n(z)+y)/3}\).
We next bound the number of levels at which a chain comes close to the
barrier.  Work first on the simultaneous barrier and oscillation event just
constructed.  Suppose that, for some boundary point $z$, the quantity
$X_i(z)+y$ visits $[0,J]$ at more than $A(J+1)^4$ levels, and let $n$ be a
level at which this bound is first exceeded.  Choose a mesh center $x_n$
within $e^{-n}$ of $z$.  At every earlier level $i\leq n$, both points lie
within a fixed mesh multiple of the same level-$i$ center.  The oscillation
bound, solved for that center's height, shows that every counted visit
gives
\[
 0\leq X_i(x_n)+y\leq C'(J+1).
\]
The lower inequality at all other levels follows from the simultaneous
barrier.  In particular the endpoint at level $n$ belongs to this same
bounded band.  This is a pathwise implication on the good event.  Having
obtained it, we drop that event and estimate the resulting event at the
fixed mesh center $x_n$.

Under the endpoint tilt by $H_n(x_n)$, put $W_i=X_i(x_n)+y$.
Its increments form the centered variance-$2$ Gaussian walk, and its
starting height is $y$ plus the bounded Gaussian correction used in
Lemma~\ref{proxy:barrier}.  Let $V_J$ count its visits to
$[0,C'(J+1)]$ before killing below zero.  This count has an exponential
upper tail on the $(J+1)^2$ scale, uniformly in the starting height.
Indeed, start the first block at the first visit to the band, give it
$L_J=\lceil C''(J+1)^2\rceil$ steps, and start each subsequent block at the
first later band visit.  From every point of the band there is a
probability at least $c>0$, independent of $J$, of being killed during that
block.
Each block contains at most $L_J$ counted visits.  The strong Markov property
therefore gives
\[
 \Prob_*(V_J\geq k)\leq C\exp\{-ck/(J+1)^2\}.
\]
At $k=A(J+1)^4$, combine this estimate with the endpoint ballot bound by
H\"older, retaining a ballot power $b>2/3$.  Undoing the tilt and summing over
the $O(e^n)$ possible centers gives the same factor $e^{-y}$ as before.
The endpoint bands cost at most $e^{C'(J+1)}$, whereas the visit estimate
contributes $e^{-cA(J+1)^2}$ after adjusting $c$.  The sums over $J$ and over
$n$, the latter bounded by
\[
 \sum_{n\geq1}n^{-3b/2}(1+y+\log(2+n))^C,
\]
are finite.  Thus, outside an event of probability
$C_d(1+y)^Ce^{-y}$, every boundary chain has at most $A(J+1)^4$ visits to
$[0,J]$, simultaneously for all integers $J\geq1$.
Let $A_*$ be the supremum of the summed cost ratios over all boundary
chains in the slab.  The preceding bounds give
\[
 A_*\leq C_d e^{dy}\sum_{J\geq0}(J+2)^4e^{-dJ/3}<\infty.
\]
The bound $\Prob(A_*>C_de^{dy})\leq C_d(1+y)^Ce^{-y}$ gives
$\E A_*^q<\infty$ for every $q<1/d$.  Since $d=(1-\theta)\xi_0$, choosing
$\theta$ sufficiently close to one supplies any prescribed moment, with
room for the H\"older inequalities used to restore the radial part.

Finitely many larger-scale slab-interior crossings to link these chains and move right one slab are handled the same way without level summation, charged against a fixed interval still to the right.
They can use common surrounding paths to connect the choices.
Enlarge $A_*$ to include these finitely many links; the same moment bounds
continue to hold.
All these estimates hold for \(P^\perp\) by the compact slab continuous Gaussian corrections.

\smallskip\noindent\textit{Restoring the radial part and the perimeter.}\quad

Restore the radial and normalization factors: per slab needed to go right from \(z\) (including the starting one) the cost is at most the right side of \eqref{eq:original-11} with \(K\) replaced by
\[
 C\,A_k V_\gamma^{-(1-\theta)2\xi_0/\gamma}
 \exp\{-d\,U_a(k)+C\,{\rm osc}_{[k-C',k+C']}U_a\}.
\]
Indeed each compared mass interval is inside the indicated tail; \(A_k\) depends only on the lateral field with the above moment control and bounded normalization-unit constants.
Sum the bounds over all \(k\): the sum of the displayed exponential occupation factors alone has all positive moments uniformly.
For integer moments order tuple times, group overlapping bounded windows (groups still bounded for fixed order), and use at each next separated group the Gaussian density bound \(C(1+m)^{-3/2}\) for the 3-dimensional drifted motion across gap \(m\), integrating the exponential radial decay with bounded-window supremum-increment exponential moments.
The two halves cause no change.
Conditional Minkowski over the lateral field therefore preserves the desired large moments of the \(A_k\)-weighted sum, with slack for Hölder with inverse \(V_\gamma\) (all orders by one slab and radial bounds), and the small positive DMS tilt.
This proves \eqref{eq:original-11}.
The same proof works at limiting charge, with limiting normalization constants.

\end{proof}

\subsection{Uniform boundary continuity and small images}\label{proxy:consequences}

Consequences are used in the following probabilistic sense.
Boundary arcs of small relative length incident to a typical boundary mark have images of small diameter, with the exceptional upper-chain probability at any fixed proxy threshold a power of relative length of arbitrarily high prescribed order, times the corresponding perimeter scale factor.
If an arc passes the second typical strip mark, use \eqref{eq:original-11} at both ends.
To treat \emph{all} arcs of a given size discretize by a mesh of that size from a uniform length origin, each anchor again marginally typical after disintegrating the cut process as necessary; use coordinates based on each anchor.
This loses only one power for the grid union.
The whole boundary can also be covered with cost scale \(s^{2\xi_0/\gamma}\), with such moments.
In these assertions the upper costs just require paths between near-boundary points in the open parent; passing through Lemma~\ref{proxy:lower-cost}, a single high-probability lower bound in the parent, gives the actual Euclidean smallness simultaneously, using continuous closure maps.
Whole-domain diameter then follows also by the maximum principle.
Equicontinuity in prime coordinates of disks with lengths in compact positive ranges follows in probability by the above and convergence of the standard-coordinate boundary measures to diffuse full-support measures; bounded analytic maps with equicontinuous boundaries are equicontinuous on the closed disk.
None of this asserts independence of coordinates of the embeddings from the cut surfaces used for \eqref{eq:original-11}.

